\documentclass[11pt,a4paper]{article}
\usepackage[margin=1in]{geometry}
\usepackage{amsmath,amssymb,amsthm}
\usepackage{booktabs,tabularx,longtable}
\usepackage{graphicx,xcolor}
\usepackage{hyperref}
\usepackage[ruled,vlined]{algorithm2e}
\usepackage{enumitem}
\usepackage{fancyhdr}
\usepackage{titlesec}

\usepackage{tcolorbox}
\usepackage{listings}
\usepackage{caption}

\definecolor{accent}{HTML}{2563EB}
\definecolor{codebg}{HTML}{F8FAFC}
\definecolor{synergy}{HTML}{16A34A}
\definecolor{redundancy}{HTML}{DC2626}

\hypersetup{colorlinks=true,linkcolor=accent,citecolor=accent,urlcolor=accent}

\lstset{
  language=Python, basicstyle=\ttfamily\small,
  backgroundcolor=\color{codebg}, frame=single,
  rulecolor=\color{gray!30}, breaklines=true,
  keywordstyle=\color{accent}\bfseries,
  commentstyle=\color{gray},
  numbers=left, numberstyle=\tiny\color{gray},
  xleftmargin=2em, framexleftmargin=1.5em
}

\tcbuselibrary{skins,breakable}
\newtcolorbox{notebox}[1][]{
  colback=accent!5, colframe=accent!60, fonttitle=\bfseries,
  title=#1, breakable, sharp corners=south
}

\pagestyle{fancy}
\fancyhf{}
\fancyhead[L]{\small\textcolor{gray}{Research Plan — Constraint-Interaction Metric}}
\fancyhead[R]{\small\textcolor{gray}{Biswajyoti Nath}
\fancyfoot[C]{\thepage}

\titleformat{\section}{\Large\bfseries\color{accent}}{\thesection}{1em}{}
\titleformat{\subsection}{\large\bfseries}{\thesubsection}{1em}{}

\title{%
  \vspace{-1cm}
  {\LARGE\bfseries\color{accent} Research Plan}\\[0.3cm]
  {\Large Constraint-Interaction Metric for Symbolic Regression}\\[0.5cm]
  {\normalsize\textit{Summer Research Project — July 6–28, 2026}}
}
\author{%
  \textbf{Biswajyoti Nath}\\
  Department of Computer Science \& Engineering\\
  Assam Science and Technology University
}
\date{}

\begin{document}
\maketitle
\thispagestyle{empty}

\begin{abstract}
We propose a \textbf{constraint-interaction metric} $M(i,j)$ for symbolic regression (SR) that quantifies whether pairs of search constraints act synergistically ($M>1$), redundantly ($M<1$), or independently ($M\approx1$). The metric is defined as $M(i,j) = \rho(C_i \wedge C_j)/[\rho(C_i)\rho(C_j)]$, where $\rho(C)$ is the fraction of grammatically valid expressions satisfying constraint $C$. We will build a prototype system integrating a grammar-based expression generator, Monte Carlo density estimation, and PySR-based experiments on Feynman physics benchmarks. Deliverables include an open-source code repository, a 55--65 page implementation handbook, and empirical evidence that $M$ predicts performance gains (target: Pearson $r > 0.6$, $p < 0.05$). Timeline: 23 working days (July 6--28, 2026).
\end{abstract}

\tableofcontents
\newpage

%% ======================================================================
\section{Goals and Scope}
\label{sec:goals}

\subsection{Primary Goals}
\begin{enumerate}[label=\textbf{G\arabic*}]
  \item Formally define $\rho(C)$ and the interaction coefficient $M(i,j)$.
  \item Implement efficient Monte Carlo algorithms to estimate these probabilities.
  \item Embed the framework in a symbolic regression workflow using PySR.
  \item Empirically validate that $M$ predicts how constraint combinations affect search performance.
\end{enumerate}

\subsection{Scope}
Given the 23-day timeline, we target a \textbf{focused prototype} producing convincing initial results suitable for a workshop or short conference submission. Specifically:
\begin{itemize}
  \item A working system: data generators, metric calculator, experiment scripts.
  \item A moderate experiment set: 3--4 datasets, 4--5 constraint types, key pairwise combinations.
  \item Target: demonstrate $M(i,j)$ correlates with performance gains (Pearson $r > 0.6$, $p < 0.05$ after FDR correction) across representative SR problems.
\end{itemize}

\subsection{Deliverables}
\begin{table}[h]
\centering\small
\begin{tabularx}{\textwidth}{lX}
\toprule
\textbf{Deliverable} & \textbf{Format} \\
\midrule
Implementation handbook & \texttt{Handbook.pdf} (55--65 pages) \\
Code repository & GitHub (Python modules, notebooks) \\
Sample datasets & \texttt{data/} folder (CSV) \\
Figures and results & \texttt{figures/}, \texttt{results/} (PNG/CSV) \\
README \& CI & \texttt{README.md}, \texttt{.github/workflows/ci.yml} \\
\bottomrule
\end{tabularx}
\caption{Project deliverables.}
\end{table}

%% ======================================================================
\section{Background and Related Work}
\label{sec:background}

Symbolic regression (SR) discovers interpretable mathematical expressions fitting data. It is NP-hard due to the combinatorial search space. Researchers use \textit{constraints}---grammar rules, type systems, operator whitelists, dimensional analysis---to incorporate domain knowledge and reduce search effort.

\subsection{Prior Work on Constraints in SR}
\begin{itemize}
  \item \textbf{Grammar-based SR:} Attribute grammars enforcing dimensional consistency recovered ${\sim}60\%$ more benchmark equations~(Brence et al., 2023).
  \item \textbf{Semantic backpropagation:} Reissmann et al.\ (2025) use semantic constraints during evolution in GEP.
  \item \textbf{Graph-based SR:} Xiang et al.\ (2025) encode constraints into neural-guided MCTS.
  \item \textbf{PySR:} State-of-the-art open-source SR using multi-population evolutionary search (Cranmer, 2023).
\end{itemize}

\subsection{Research Gap}
\begin{notebox}[Novelty Statement]
All prior work studies individual constraint effectiveness. \textbf{No existing work models pairwise statistical interaction between constraints as a predictive signal for search performance.} Our metric $M(i,j)$ fills this gap. While analogous to ``lift'' in association rule mining, the application to structured expression grammars---where co-occurrence is measured over a \textit{generative} distribution rather than an observed corpus---is novel in the SR literature.
\end{notebox}

%% ======================================================================
\section{Proposed Metric}
\label{sec:metric}

\subsection{Definitions}
Let $\mathcal{E}$ be the space of expressions generated by grammar $G$. For a constraint $C: \mathcal{E} \to \{0,1\}$, define:
\begin{equation}
  \rho(C) \;=\; \Pr_{E \sim G}[C(E) = 1]
\end{equation}

The \textbf{interaction coefficient} for constraints $C_i, C_j$ is:
\begin{equation}
  \boxed{M(i,j) \;=\; \frac{\rho(C_i \wedge C_j)}{\rho(C_i)\,\rho(C_j)}}
  \label{eq:metric}
\end{equation}

\subsection{Interpretation}
\begin{itemize}
  \item $M(i,j) > 1$: \textcolor{synergy}{\textbf{Synergy}} --- constraints co-occur more than chance; combining them may yield super-additive gains.
  \item $M(i,j) < 1$: \textcolor{redundancy}{\textbf{Redundancy}} --- one constraint subsumes the other; adding the second provides diminishing returns.
  \item $M(i,j) = 1$: \textbf{Independence} --- constraints act on orthogonal dimensions.
\end{itemize}

\subsection{Statistical Significance of Interaction}
Rather than using arbitrary thresholds ($M > 1.1$, $M < 0.9$), we determine significance via a bootstrap confidence interval. Under the null hypothesis of independence, $M = 1$. We:
\begin{enumerate}
  \item Bootstrap-resample the expression set $B = 1000$ times.
  \item Compute $\hat{M}^{(b)}$ for each resample.
  \item Construct a 95\% CI for $M$. If $1 \notin \text{CI}$, reject independence.
\end{enumerate}

%% ======================================================================
\section{Data Generation}
\label{sec:datagen}

\subsection{Grammar-Based Expression Generator}
We use a context-free grammar in BNF:
\begin{lstlisting}[language={},frame=none,backgroundcolor=\color{white}]
<expr> ::= <expr> <binop> <expr>
         | <unop>(<expr>)
         | <variable> | <constant>
<binop> ::= + | - | *
<unop>  ::= sin | cos
<variable> ::= x | y
<constant> ::= 1 | 2 | ... | 9
\end{lstlisting}

A recursive generator applies grammar rules up to \texttt{max\_depth}. Division is excluded from the default grammar to avoid domain errors; it can be added as a configurable extension.

\subsection{Sampling Strategy}
\begin{table}[h]
\centering\small
\begin{tabular}{lll}
\toprule
\textbf{Parameter} & \textbf{Value} & \textbf{Rationale} \\
\midrule
Grammar depth limit & 8 & Moderate complexity \\
Operators & $+, -, \times, \sin, \cos$ & Core arithmetic + trig \\
Variables & $x, y$ & Two-input functions \\
Constants & $\pm 1 \ldots 9$ & Discrete values \\
Sample size $N$ & 100{,}000 & $\rho \sim 0.01 \Rightarrow {\sim}1000$ hits, SE ${\sim}3\%$ \\
Stratification & Equal samples per depth 2,4,6,8 & Avoid depth bias \\
\bottomrule
\end{tabular}
\caption{Monte Carlo sampling parameters.}
\end{table}

Each expression string is converted to a SymPy object via \texttt{sympify()}, with invalid expressions (syntax errors, domain violations) discarded and replaced.

%% ======================================================================
\section{Constraint Definitions}
\label{sec:constraints}

We implement five constraints $C(E)$, designed to be \textbf{structurally distinct}:

\begin{enumerate}[label=\textbf{C\arabic*}]
  \item \textbf{Structural Template ($C_{\mathrm{struct}}$):} The expression must conform to a structural pattern---e.g., no nested trigonometric calls, at most 2 consecutive binary operations. This is checked by traversing the SymPy AST.

  \item \textbf{Maximum Depth ($C_{\mathrm{depth}}$):} Parse-tree depth $\le$ threshold (default: 6). Recursive depth computation on \texttt{sym\_expr}.

  \item \textbf{Operator Whitelist ($C_{\mathrm{op}}$):} Only operators from a prescribed set (e.g., $\{+, -, \times\}$) may appear. Every AST node is verified against the whitelist.

  \item \textbf{Positivity ($C_{\mathrm{pos}}$):} The expression evaluates to $\ge 0$ for all inputs in $[-10, 10]^2$. We test on \textbf{200 random points} (not 10) and report the empirical false-positive rate: $P(\text{pass} \mid \text{violates}) \le (0.95)^{200} < 10^{-4}$ if 5\% of the domain violates.

  \item \textbf{Dimensional Consistency ($C_{\mathrm{dim}}$, optional):} Each variable carries a physical dimension; we propagate units through operations and reject inconsistencies (e.g., $\text{length} + \text{time}$). Skipped if time is short.
\end{enumerate}

\begin{notebox}[Design Decision: C1 vs.\ C3]
$C_{\mathrm{struct}}$ and $C_{\mathrm{op}}$ are intentionally distinct: the former constrains \textit{tree shape} (nesting patterns), the latter constrains \textit{node labels} (which operators appear). Their $M$ value is therefore non-trivial.
\end{notebox}

%% ======================================================================
\section{Monte Carlo Estimation}
\label{sec:montecarlo}

\begin{algorithm}[H]
\SetAlgoLined
\KwIn{Grammar $G$, constraints $\{C_1, \ldots, C_k\}$, sample size $N$}
\KwOut{Densities $\hat{\rho}(C_i)$, joint densities $\hat{\rho}(C_i \wedge C_j)$, matrix $M$}
Initialize counters $n_i \gets 0$, $n_{ij} \gets 0$\;
\For{$t = 1$ \KwTo $N$}{
  $E \gets \texttt{generate\_expr}(G, \texttt{max\_depth}=8)$\;
  $s \gets \texttt{sympify}(E)$\;
  \If{$s$ is invalid}{\textbf{continue} (replace sample)\;}
  \For{$i = 1$ \KwTo $k$}{
    \If{$C_i(s)$}{$n_i \gets n_i + 1$\;}
    \For{$j = i+1$ \KwTo $k$}{
      \If{$C_i(s) \wedge C_j(s)$}{$n_{ij} \gets n_{ij} + 1$\;}
    }
  }
}
$\hat{\rho}(C_i) \gets n_i / N$\;
$\hat{\rho}(C_i \wedge C_j) \gets n_{ij} / N$\;
$M(i,j) \gets \hat{\rho}(C_i \wedge C_j) / [\hat{\rho}(C_i) \cdot \hat{\rho}(C_j)]$\;
\Return{$\hat{\rho}, M$}
\caption{Monte Carlo Constraint Density Estimation}
\label{alg:montecarlo}
\end{algorithm}

\textbf{Parallelism:} Sampling is embarrassingly parallel. We use Python's \texttt{multiprocessing} or \texttt{joblib} to distribute across cores.

\textbf{Convergence check:} Run at $N = 1\text{k}, 10\text{k}, 100\text{k}$ and verify $|\hat{\rho}_{100\text{k}} - \hat{\rho}_{10\text{k}}| < 0.01$.

%% ======================================================================
\section{Software Architecture}
\label{sec:architecture}

\subsection{Module Structure}
\begin{table}[h]
\centering\small
\begin{tabularx}{\textwidth}{lX}
\toprule
\textbf{Module} & \textbf{Responsibility} \\
\midrule
\texttt{expr\_generator.py} & Grammar-based recursive expression generation \\
\texttt{constraints.py} & Constraint implementations (C1--C5) with unit tests \\
\texttt{metric.py} & Monte Carlo sampling, $\rho$ estimation, $M$ computation \\
\texttt{experiments.py} & PySR integration, experiment orchestration, logging \\
\texttt{analysis.py} & Statistical analysis, correlation, regression, plotting \\
\texttt{config.yaml} & All parameters (grammar, constraints, PySR, seeds) \\
\texttt{main.py} & CLI entry point: \texttt{-{}-phase=\{metric,experiment,analysis\}} \\
\bottomrule
\end{tabularx}
\caption{Module structure.}
\end{table}

\subsection{Configuration Schema}
\begin{lstlisting}[caption={Minimal \texttt{config.yaml} schema.}]
grammar:
  max_depth: 8
  operators: ["+", "-", "*", "sin", "cos"]
  variables: ["x", "y"]
monte_carlo:
  N: 100000
  stratify_by_depth: true
pysr:
  population_size: 40
  niterations: 100
  timeout_seconds: 300
datasets: ["feynman_1", "feynman_2", "polynomial"]
seeds: [42, 123, 456, 789, 1024]
\end{lstlisting}

\subsection{Testing Strategy}
\begin{itemize}
  \item \textbf{Unit tests:} Known pass/fail expressions for each constraint.
  \item \textbf{Convergence test:} Assert $\rho$ stabilizes within 5\% between $N=10\text{k}$ and $N=100\text{k}$.
  \item \textbf{Integration test:} End-to-end pipeline on a toy 3-operator grammar.
\end{itemize}

\subsection{Dependency Management}
All dependencies pinned in \texttt{requirements.txt}:
\texttt{pysr>=0.19}, \texttt{sympy==1.12}, \texttt{numpy>=1.24}, \texttt{scipy>=1.11}, \texttt{seaborn>=0.13}, \texttt{pandas>=2.0}.

%% ======================================================================
\section{Experimental Design}
\label{sec:experiments}

\subsection{Datasets}
\begin{enumerate}[label=\textbf{D\arabic*}]
  \item \textbf{Feynman-1:} Simple physics law (e.g., $E = mc^2$, kinetic energy).
  \item \textbf{Feynman-2:} Multi-variable with trigonometry (e.g., projectile motion).
  \item \textbf{Polynomial:} Synthetic 3-term polynomial.
  \item \textbf{SRSD:} Dataset with dummy variable (Matsubara et al., 2024).
\end{enumerate}

\subsection{Constraint Scenarios}
For each dataset, PySR is run under:
\begin{itemize}
  \item \textbf{Baseline:} No additional constraints.
  \item \textbf{Singles:} Each of C1--C4 individually.
  \item \textbf{Key pairs:} Pairs with extreme predicted $M$ values (highest synergy, highest redundancy).
  \item \textbf{All constraints:} C1--C4 combined.
\end{itemize}
Each scenario uses \textbf{5 repeats} (fixed seeds) for statistical reliability.

\subsection{Constraint Enforcement in PySR}
Constraints are enforced \textbf{during search} via PySR's \texttt{constraints} and \texttt{nested\_constraints} parameters (for structural/operator constraints) and custom loss penalties (for positivity). This ensures the metric captures genuine search-effort reduction, not post-hoc filtering.

\subsection{Success Criterion}
Recovery is defined as: NED $< 0.1$ \textbf{or} relative MSE $< 10^{-6}$ on 20\% held-out data.

\subsection{Metrics Recorded}
\begin{itemize}
  \item Wall-clock runtime (seconds)
  \item Number of candidate evaluations
  \item Recovery rate (binary)
  \item Solution complexity (AST node count)
\end{itemize}

%% ======================================================================
\section{Statistical Analysis}
\label{sec:stats}

\subsection{Speedup Definition}
For constraint pair $(i,j)$, the \textbf{speedup} compares combined performance against the best single:
\begin{equation}
  \boxed{S(i,j) = \frac{\min(\text{evals}_{C_i},\; \text{evals}_{C_j})}{\text{evals}_{C_i \wedge C_j}}}
  \label{eq:speedup}
\end{equation}
$S > 1$ means the combination outperforms either constraint alone.

\subsection{Correlation Analysis}
\begin{itemize}
  \item Pearson and Spearman correlation between $M(i,j)$ and $S(i,j)$.
  \item Linear regression: $S(i,j) = a + b \cdot M(i,j) + \varepsilon$. Significant positive $b$ ($p < 0.05$) supports predictive value.
  \item \textbf{Multiple comparisons:} Benjamini--Hochberg FDR correction across all pairwise tests.
\end{itemize}

\subsection{Confidence Intervals}
95\% bootstrap CIs for all mean metrics. Error bars on all plots.

\subsection{Power Considerations}
With $\binom{4}{2} \times 4 = 24$ (pair, dataset) data points, a power analysis (\texttt{statsmodels.stats.power}) confirms detectability of $r \ge 0.6$ at $\alpha = 0.05$ with power $> 0.80$.

%% ======================================================================
\section{Visualization Plan}
\label{sec:viz}
\begin{enumerate}
  \item \textbf{Interaction Heatmap:} $M(i,j)$ matrix with \texttt{seaborn.heatmap}, annotated cells, coolwarm colormap.
  \item \textbf{Scatter Plot:} $M(i,j)$ vs.\ $S(i,j)$ with regression line and $r$/$p$ annotation.
  \item \textbf{Box Plots:} Runtime/evaluations per constraint scenario per dataset.
  \item \textbf{Convergence Plot:} $\hat{\rho}(C)$ vs.\ sample size $N$.
\end{enumerate}
All generated with \texttt{matplotlib}/\texttt{seaborn}; publication-quality (300 DPI, vector where possible).

%% ======================================================================
\section{Potential Pitfalls and Mitigations}
\label{sec:pitfalls}

\begin{table}[h]
\centering\small
\begin{tabularx}{\textwidth}{lXX}
\toprule
\textbf{Risk} & \textbf{Impact} & \textbf{Mitigation} \\
\midrule
Sampling noise for rare $\rho$ & Inaccurate $M$ & Increase $N$; report bootstrap CIs \\
Generator $\neq$ search distribution & $M$ may not transfer & Compare with PySR gen-0 population \\
SymPy edge cases & Constraint bugs & Unit tests on known expressions \\
PySR stochasticity & Wide CIs & 5 repeats; fix seeds \\
PySR installation failure & Blocked experiments & Fallback to \texttt{gplearn} \\
Timeline slippage & Incomplete deliverables & 2 buffer days; cut $C_{\mathrm{dim}}$ first \\
\bottomrule
\end{tabularx}
\caption{Risk register.}
\end{table}

%% ======================================================================
\section{Timeline}
\label{sec:timeline}

\begin{longtable}{clp{9cm}c}
\toprule
\textbf{Day} & \textbf{Date} & \textbf{Tasks} & \textbf{Hours} \\
\midrule
\endfirsthead
\toprule
\textbf{Day} & \textbf{Date} & \textbf{Tasks} & \textbf{Hours} \\
\midrule
\endhead
1 & Jul 6 & Setup env; install Julia+PySR; review literature; begin generator & 6 \\
2 & Jul 7 & Finish generator; SymPy integration; depth-limit tests & 8 \\
3 & Jul 8 & Monte Carlo sampling; validity checks; benchmark 1k samples & 8 \\
4 & Jul 9 & Implement C1 (structural), C2 (depth), C3 (operator) & 7 \\
5 & Jul 10 & Implement C4 (positivity, 200 pts); optional C5 (dimensional) & 8 \\
6 & Jul 11 & Full Monte Carlo run ($N{=}100\text{k}$); convergence check & 8 \\
7 & Jul 12 & Compute $M$ matrix; heatmap visualization; identify patterns & 6 \\
8 & Jul 13 & Document theory ($\rho$, $M$, bootstrap CI); write handbook \S1--4 & 7 \\
9 & Jul 14 & Software architecture; module APIs; pseudocode in handbook & 6 \\
10 & Jul 15 & PySR setup; prepare benchmark datasets; data loading scripts & 6 \\
11 & Jul 16 & Integrate constraints with PySR API; plan enforcement strategy & 7 \\
12 & Jul 17 & Baseline SR experiments (no constraints); debug; store results & 8 \\
13 & Jul 18 & Single-constraint experiments (5 repeats each); log metrics & 8 \\
14 & Jul 19 & Pair and all-constraint experiments; focus on extreme $M$ pairs & 8 \\
15 & Jul 20 & Statistical analysis: means, CIs, correlations, FDR correction & 6 \\
16 & Jul 21 & Ablation studies; begin drafting Results chapter & 6 \\
17 & Jul 22 & Create all figures (heatmap, scatter, boxplots) & 7 \\
18 & Jul 23 & Write Experiments \& Results chapters & 8 \\
19 & Jul 24 & Write Implementation Details chapter & 8 \\
20 & Jul 25 & Write Data Generation \& Constraints chapters & 7 \\
21 & Jul 26 & \textit{Buffer day:} catch-up on any delayed tasks & 6 \\
22 & Jul 27 & Finalize writing; proofread; compile handbook PDF & 6 \\
23 & Jul 28 & \textit{Buffer + submit:} final polish, push repo, CI, README & 4 \\
\bottomrule
\caption{Daily timeline (July 6--28, 2026). Days 21 and 23 are buffer days.}
\end{longtable}

%% ======================================================================
\section{Handbook Outline}
\label{sec:handbook}

Target: \textbf{55--65 pages}, ${\sim}25{,}000$ words.

\begin{enumerate}
  \item Cover Page (0.5 p.)
  \item Abstract \& Executive Summary (1 p.)
  \item Table of Contents (1 p.)
  \item Ch.\ 1: Introduction (2 p.) --- SR background, motivation, contributions.
  \item Ch.\ 2: Related Work (3 p.) --- SR techniques, constraints, benchmarks.
  \item Ch.\ 3: Problem Definition (2 p.) --- Formal $M(i,j)$, research questions.
  \item Ch.\ 4: Proposed Metric (4 p.) --- Derivation, interpretation, bootstrap CI.
  \item Ch.\ 5: Software Architecture (5 p.) --- Modules, data flow, config schema.
  \item Ch.\ 6: Data \& Expression Generation (4 p.) --- Grammar, sampling, stratification.
  \item Ch.\ 7: Constraint Definitions (4 p.) --- Formal + implementation.
  \item Ch.\ 8: Monte Carlo Estimation (3 p.) --- Algorithm, convergence, parallelism.
  \item Ch.\ 9: PySR Integration (4 p.) --- API usage, constraint enforcement.
  \item Ch.\ 10: Experimental Protocol (4 p.) --- Datasets, experiment matrix, metrics.
  \item Ch.\ 11: Statistical Methods (3 p.) --- Correlation, regression, FDR, power.
  \item Ch.\ 12: Results (5 p.) --- Heatmap, scatter plots, tables.
  \item Ch.\ 13: Discussion (4 p.) --- Interpretation, limitations.
  \item Ch.\ 14: Conclusion \& Future Work (2 p.)
  \item Appendices (4 p.) --- Code listings, extended tables.
\end{enumerate}

%% ======================================================================
\section*{References}

\begin{enumerate}[label={[\arabic*]},leftmargin=*]
  \item M.\ Reissmann et al., ``Constraining Genetic Symbolic Regression via Semantic Backpropagation,'' \textit{Genetic Programming and Evolvable Machines}, 2025.
  \item Y.\ Xiang et al., ``Graph-based Symbolic Regression with Invariance and Constraint Encoding,'' \textit{NeurIPS}, 2025.
  \item M.\ Cranmer, ``Interpretable Machine Learning for Science with PySR and SymbolicRegression.jl,'' \textit{arXiv:2305.01582}, 2023.
  \item J.\ Brence et al., ``Dimensionally-consistent equation discovery through probabilistic attribute grammars,'' \textit{Information Sciences}, vol.\ 632, 2023.
  \item Y.\ Matsubara et al., ``Rethinking Symbolic Regression Datasets and Benchmarks for Scientific Discovery,'' \textit{J.\ Data-centric Machine Learning Research}, 2024.
  \item J.\ Martin and M.\ Burnett, ``Shape-Constrained Symbolic Regression---Improving Extrapolation,'' in \textit{GECCO}, 2014.
  \item PySR Documentation and GitHub, 2023. \url{https://ai.damtp.cam.ac.uk/pysr/}
\end{enumerate}

\end{document}
